\section{Overview}
\label{sec:approach-overview}

The system takes a pull request and produces a written review. Three design
properties make the comparison in Chapter~\ref{chap:edas} a comparison of
context rather than of pipelines.

The pipeline splits into a construction stage and a generation stage, with a
frozen artefact between them. Construction reads the repository once, at the
pull request's head commit, and writes an \emph{evidence pack}: a JSON file
holding the change, the structural facts, and the retrieved code snippets
(Section~\ref{sec:approach-packs}). Generation reads that file and never
touches the repository. Because all four conditions read the same pack, none
can differ by accident of timing.

The conditions are nested. \baseline{} receives the change alone. \kg{} adds a
structural block, \rag{} adds a retrieval block, and \hybrid{} adds both. No
condition is deprived of the diff. The context arms add evidence and a sentence
that names its type, so the contrast includes both supplied information and
context-specific instruction.

Structural context is produced by a code property graph (CPG) built with
Joern, a static-analysis platform that parses source into an interlinked
representation of syntax, control flow, and data flow
(Section~\ref{sec:approach-kg}). Throughout the thesis, \kg{} and ``the
graph arm'' name this same condition. The human study uses the same
knowledge-graph block, with dependent edges resolved from Python import
statements (Section~\ref{sec:approach-kg-ast}). Figure~\ref{fig:pipeline}
shows the arrangement.

\begin{figure}[htbp]
    \centering
    \resizebox{0.92\textwidth}{!}{%
    % !TeX root = ../thesis.tex
% Pipeline figure — clean academic style (thin lines, 2 accent colours).
\begin{tikzpicture}[
    >=Latex,
    every node/.style={font=\footnotesize, inner sep=0pt},
    box/.style={
      draw=black!50, line width=0.5pt, rounded corners=2pt,
      fill=white, minimum height=7mm, text centered,
      inner sep=4pt, align=center
    },
    kgbox/.style={box, draw=kgblue!70, fill=kgblue!6},
    ragbox/.style={box, draw=ragorange!70, fill=ragorange!6},
    accentbox/.style={box, line width=0.8pt},
    arr/.style={->, line width=0.5pt, black!50},
    fatarr/.style={->, line width=0.7pt, black!55}
]

% ── Row 0: Repository ─────────────────────────────────
\node[box, minimum width=6cm] (repo)
    at (5.5, 0) {Repository at PR head commit};

% ── Row 1: Three extractors ───────────────────────────
\node[box, minimum width=2.4cm] (diff) at (1.2,-1.4) {Unified diff};
\node[kgbox, minimum width=2.8cm] (cpg)  at (5.5,-1.4) {Joern CPG};
\node[ragbox, minimum width=2.8cm] (emb)  at (9.8,-1.4) {Embedding index};

\draw[arr] (repo) -- (diff);
\draw[arr] (repo) -- (cpg);
\draw[arr] (repo) -- (emb);

% ── Row 2: KG detail ─────────────────────────────────
\node[kgbox, minimum width=2.0cm, font=\scriptsize] (calls)
    at (3.7,-2.8) {Callers};
\node[kgbox, minimum width=2.0cm, font=\scriptsize] (deps)
    at (5.9,-2.8) {Dependents};
\node[box, minimum width=2.0cm, font=\scriptsize] (tests)
    at (8.1,-2.8) {Tests\\[-2pt]{\tiny filename match}};

\draw[arr] (cpg) -- (calls);
\draw[arr] (cpg) -- (deps);
\draw[arr] ([xshift=2.6cm]repo.south) -- (tests.north);

% ── Row 2: RAG detail ─────────────────────────────────
\node[ragbox, minimum width=2.4cm, font=\scriptsize] (chunks)
    at (10.5,-2.8) {Similar chunks};
\draw[arr] (emb) -- (chunks);

% ── Row 3: Evidence pack ──────────────────────────────
\node[accentbox, minimum width=10.5cm, minimum height=9mm,
      draw=black!60, fill=black!3] (pack)
    at (5.5,-4.2) {\textbf{Evidence pack}\;%
    {\scriptsize (one JSON file per PR)}};

\draw[arr] (diff.south) -- (diff.south |- pack.north);
\draw[arr] (calls.south) -- (calls.south |- pack.north);
\draw[arr] (deps.south) -- (deps.south |- pack.north);
\draw[arr] (tests.south) -- (tests.south |- pack.north);
\draw[arr] (chunks.south) -- (chunks.south |- pack.north);

% ── Annotation: "generated once, read by all modes" ───
\node[font=\scriptsize\itshape, text=black!45, anchor=west]
    at (11.0,-4.2) {generated once};

% ── Row 4: Four modes ─────────────────────────────────
\node[box, minimum width=2.1cm, fill=black!3] (bl)
    at (1.85,-5.8) {\textbf{Baseline}\\[-1pt]{\scriptsize diff}};
\node[kgbox, minimum width=2.1cm] (kgm)
    at (4.55,-5.8) {\textbf{KG}\\[-1pt]{\scriptsize diff\,+\,graph}};
\node[ragbox, minimum width=2.1cm] (ragm)
    at (7.25,-5.8) {\textbf{RAG}\\[-1pt]{\scriptsize diff\,+\,chunks}};
\node[box, minimum width=2.1cm, draw=black!50,
      fill=black!2] (hyb)
    at (9.95,-5.8) {\textbf{Hybrid}\\[-1pt]{\scriptsize diff\,+\,both}};
% Hybrid dual-colour accent
\draw[kgblue!50, line width=0.6pt, rounded corners=2.5pt]
    (8.84,-5.28) rectangle (11.06,-6.32);
\draw[ragorange!50, line width=0.6pt, rounded corners=3pt]
    (8.78,-5.22) rectangle (11.12,-6.38);

\draw[arr] (pack.south) -- ++(0,-0.35) -| (bl.north);
\draw[arr] (pack.south) -- ++(0,-0.35) -| (kgm.north);
\draw[arr] (pack.south) -- ++(0,-0.35) -| (ragm.north);
\draw[arr] (pack.south) -- ++(0,-0.35) -| (hyb.north);

% ── Row 5: Generator ─────────────────────────────────
\node[box, minimum width=9.5cm] (gen)
    at (5.5,-7.2) {GPT-4o generator \;\;$\rightarrow$\;\;
    {\scriptsize $35 \times 4 = 140$ reviews}};

\draw[arr] (bl.south)  -- (bl.south |- gen.north);
\draw[arr] (kgm.south) -- (kgm.south |- gen.north);
\draw[arr] (ragm.south) -- (ragm.south |- gen.north);
\draw[arr] (hyb.south) -- (hyb.south |- gen.north);

% ── Row 6: Judge panel ────────────────────────────────
\node[box, minimum width=10.5cm, minimum height=9mm] (judge)
    at (5.5,-8.6)
    {5-model judge panel\;\;$\times$\;\;25 criteria\;\;$=$\;\;%
     {\scriptsize 17\,500 scored cells}};

\draw[fatarr] (gen) -- (judge);

% ── Bracket + judge names on right ────────────────────
\draw[black!30, line width=0.4pt]
    (11.0,-7.9) -- (11.2,-7.9) -- (11.2,-9.3) -- (11.0,-9.3);
\node[font=\scriptsize, text=black!45, anchor=west, align=left]
    at (11.4,-8.6)
    {GPT-4o\\[-1pt]Gemini\,2.5\,Flash\\[-1pt]Claude\,Sonnet\,4.5\\[-1pt]%
     DeepSeek\,v4\,Pro\\[-1pt]Grok\,4.6};

\end{tikzpicture}
%
    }
    \caption{End-to-end pipeline. Arrows denote data flow: each stage
             reads the output of the stage above it. The repository is
             processed once into a frozen evidence pack. Each mode reads the
             same pack and differs only in which fields are rendered into the
             prompt. Related tests are matched by filename, not resolved from
             the graph. A five-model judge panel scores every review on 25
             binary criteria. The human study uses the same pipeline with the
             graph's dependent edges resolved from import statements
             (Section~\ref{sec:approach-kg-ast}).}
    \label{fig:pipeline}
\end{figure}

\section{Evidence Packs}
\label{sec:approach-packs}

An evidence pack is the unit of input to generation. It contains the pull
request's metadata, the list of changed files, the unified diff, the
structural facts from the graph builder, and the retrieved code chunks.

\section{Knowledge Graph Construction}
\label{sec:approach-kg}

The structural context supplied to the reviewer is a subgraph of the
repository's code property graph (CPG), built with Joern, a
language-agnostic static-analysis platform~\cite{yamaguchi2014cpg}. The CPG
represents syntax, control flow, and data flow in a single queryable
structure; the pipeline queries it for the cross-file relationships that a
diff-only reviewer cannot see.

\subsection{The Code Property Graph}
\label{sec:approach-kg-cpg}

At the pull request's head commit the full repository is loaded into Joern,
which parses every source file into an in-memory code property graph. The
pipeline then queries the graph for five relationship types seeded from the
changed files:
%
\begin{enumerate}[leftmargin=*,itemsep=2pt,label=(\roman*)]
  \item incoming callers of changed functions, with caller name, file, and
        line number;
  \item outgoing callees;
  \item file-level dependents (files that import a changed module);
  \item related test files, identified by filename-convention matching; and
  \item code owners read from the repository's \texttt{CODEOWNERS} file when
        present.
\end{enumerate}
%
Queries use exact name matching against the method and type nodes that Joern
extracts: a changed function's fully qualified name is looked up in the
call-graph overlay, and every node whose call target matches is returned.

Same-file and duplicate edges are removed because the diff already exposes
same-file context; the renderer presents up to twelve edges, ordered by
query type, without learned relevance ranking.
% Sources: experiments/2026-09-19_joern_replace_grep/,
% prnote/joern_kg.py, prnote/kg.py.

Experiment~2 additionally decomposes the graph arm into component
conditions: one containing only file-level dependency pointers, one
containing only the resolved call-graph edges, and one that adds
independently resolved inheritance relations. This design isolates the
contribution of call-graph edges from that of coarser pointers and measures
the gain from inheritance edges that the base Joern queries omit.
% Sources: experiments/2026-07-05_injection_exp2/harness/run_modes.py.

\subsection{The Knowledge-Graph Block in the Human Study}
\label{sec:approach-kg-ast}

The human study uses the same knowledge-graph context block as
Experiment~1: the same relations (changed files, dependents, related tests)
in the same rendered format. For its Python stimuli the dependent edges are
resolved from \texttt{import} statements rather than from the Joern CPG:
each candidate file is parsed with Python's standard library, its
\texttt{import} and \texttt{from} statements collected, and an edge
recorded only if a module name matches the changed file's module path. These edges are correct but file-level, and a
module-level import does not establish that a change to one function inside
that module reaches the importer.

\subsection{Rendering}
\label{sec:approach-kg-render}

Regardless of the builder, the pack's structural fields are rendered into a
Markdown block headed \emph{Knowledge Graph Context}, with sections for
changed files, related tests, dependents, code owners, functions, and callers.
The dependents, callers, and functions come from the CPG queries described
above or, in the human study, from import statements. Related tests are identified by filename-convention
matching (files whose path contains \texttt{test} and that share a directory
ancestor with a changed file), and code owners are read from the repository's
\texttt{CODEOWNERS} file when one exists. These two fields are therefore
text-matched, not statically resolved. Both constructions share the same
renderer and evidence-pack schema, so they differ in what they put in, not
how it is presented to the model.

Figure~\ref{fig:kg-block-example} shows a truncated graph block as the model
receives it.

\begin{figure}[htbp]
\centering
\begin{minipage}{0.92\textwidth}
{\footnotesize
\begin{verbatim}
**Knowledge Graph Context:**
**Changed Files:** .../Drawer.tsx, .../Drawer.story.tsx,
  .../SaveDashboardDrawer.tsx, .../InspectContent.tsx  [7 files]
**Related Tests:** .../RolePickerDrawer.test.tsx,
  .../SaveDashboardDrawer.test.tsx,
  .../HelpWizard.test.tsx                              [8 files]
**Files that depend on changes:**
  .../SaveDashboardDrawer.tsx,
  .../InspectContent.tsx                               [3 files]
**Functions in Changed Files:**
  SaveDashboardDrawerComponent (SaveDashboardDrawer.tsx);
  InspectContent (InspectContent.tsx);
  HelpWizard (HelpWizard.tsx)                         [15 entries]
**Call-graph edges (caller -> callee):**
  SaveDashboardDrawerComponent -> renderBody;
  InspectContent -> getErrors;
  InspectContent -> formatStats                        [4 edges]
\end{verbatim}
}
\end{minipage}
\caption{Truncated graph block from Grafana PR~\#96268 (Joern CPG). Paths
         are abbreviated; bracket counts show the full list lengths. The model
         receives this block between the diff and the closing instruction.}
\label{fig:kg-block-example}
\end{figure}

\section{Retrieval-Augmented Context}
\label{sec:approach-rag}

The retrieval arm is conventional dense retrieval, included as the comparison
that makes the graph arm interpretable. The repository is walked at the same
commit, skipping vendored directories and keeping recognised source
extensions; the first 200 candidate files are chunked and embedded with
\texttt{text-embedding-3-small}. Each changed file is embedded as a query,
and the five highest-scoring chunks from outside the changed files are
rendered into the prompt with their path, line range and a code excerpt.

\section{The Four Review-Generation Modes}
\label{sec:approach-modes}

A mode is a pair: a system prompt and a context builder.

\begin{table}[htbp]
    \centering
    \caption{The four generation modes. Every mode receives the title,
             description and diff; they differ in the context block appended.}
    \label{tab:modes}
    \small
    \begin{tabular}{@{}ll@{}}
        \toprule
        Mode & Context block \\
        \midrule
        \baseline{} & none \\
        \kg{}       & graph facts (dependents, callers, tests) \\
        \rag{}      & five retrieved snippets \\
        \hybrid{}   & both, concatenated \\
        \bottomrule
    \end{tabular}
\end{table}

The \hybrid{} mode is a concatenation, not a fusion. Its block is the graph
block followed by the retrieval block, with no joint ranking and no shared
budget. Its result therefore evaluates this simple concatenation, not all
possible ways of combining structural and similarity context.
Appendix~\ref{app:review-example-chapter} reproduces a side-by-side example
of a baseline and graph-augmented review for one pull request.

\section{Prompt Design}
\label{sec:approach-prompts}

All four modes share one system prompt, which casts the model as a senior
engineer and prescribes a five-heading output format. Each mode appends a
single paragraph naming the kind of context it has been given. The prompts are
reproduced in Appendix~\ref{app:prompts}. Two properties are load-bearing: the
fixed output format constrains what the evaluation can detect, and the per-arm
paragraph means the arms differ by one sentence of instruction as well as by
the block itself. The user message is assembled in a fixed order: title,
description, diff, context block, closing instruction.
